\documentclass[11pt]{article}

\usepackage{kotex}
\usepackage{amsmath, amssymb, amsthm}
\usepackage{hyperref}
\usepackage{geometry}
\geometry{margin=1in}

\theoremstyle{definition}
\newtheorem{definition}{Definition}[section]
\newtheorem{remark}[definition]{Remark}
\theoremstyle{plain}
\newtheorem{theorem}[definition]{Theorem}
\newtheorem{lemma}[definition]{Lemma}
\newtheorem{corollary}[definition]{Corollary}

\newcommand{\MSO}{\mathrm{MSO}}
\newcommand{\MSOtwo}{\mathrm{MSO}_2}
\newcommand{\MSOone}{\mathrm{MSO}_1}
\newcommand{\lean}[1]{\texttt{#1}}
\newcommand{\leanfile}[1]{\texttt{#1}}

\title{GraphMSO Lean 형식화 해설 I: Basics (Syntax and Semantics)}
\author{}
\date{}

\begin{document}
\maketitle

\tableofcontents

\section{이 문서 묶음에 대하여}

이 문서는 Lean 4 프로젝트 \leanfile{GraphMSO}의 코드를 Lean syntax를 모르는
독자가 읽을 수 있도록 옮긴 해설서 다섯 편 중 첫 번째다.
다섯 편의 역할은 다음과 같다.

\begin{itemize}
  \item \leanfile{Lean\_Basics.tex} (이 문서) --- 여러 곳에서 공통으로 쓰이는
    syntax와 semantics, 그리고 tree decomposition의 기본 API.
  \item \leanfile{Lean\_Nice.tex} --- 주어진 tree decomposition으로부터
    nice tree decomposition을 만들어내는 algorithm.
    참조: \leanfile{Courcelle/nice\_tree\_decomp.tex}.
  \item \leanfile{Lean\_MSO.tex} --- graph 위의 $\MSOtwo$ formula를
    $\tau_P$-graph 위의 $\MSO$ formula로, 다시 tree 위의 tree language로
    옮기는 translation.
    참조: \leanfile{Courcelle/lecture\_note\_expanded.tex}.
  \item \leanfile{Lean\_Automata.tex} --- tree automata API.
    참조: \leanfile{Courcelle/Thatcher-Wright-expanding.tex}.
  \item \leanfile{Lean\_Algorithm.tex} --- 실제로 실행 가능한 decision algorithm.
\end{itemize}

다른 문서의 항목을 인용할 때는 \emph{파일 이름과 label}을 함께 적는다.
예: ``\leanfile{Lean\_MSO.tex}의 \lean{mso:def:encode}''.

\subsection{읽기 전에: 최소한의 Lean syntax}
\label{basics:sec:leansyntax}

코드를 인용할 때 반복적으로 등장하는 표기만 정리한다. mathlib에 이미
있는 개념(\lean{SimpleGraph}, \lean{Walk}, \lean{Finset}, \lean{Set.ncard}
등)은 따로 설명하지 않는다.

\begin{itemize}
  \item \lean{inductive} --- inductive type을 정의한다. constructor의 목록을
    주면 그 constructor들로 생성되는 가장 작은 집합이 정의된다. formula나
    term처럼 ``syntax''를 정의할 때 쓴다.
  \item \lean{structure} --- 여러 개의 데이터(field)를 묶은 tuple이다.
    수학에서 ``집합 $V$, 그 위의 graph $G$, 그리고 predicate들의 family
    $\mathrm{pred}$의 삼중쌍''이라고 쓰는 것과 같다. field 중에는 \emph{axiom}
    (증명해야 할 proposition)이 들어갈 수도 있다. 이 경우 그 structure의
    원소를 만들려면 axiom의 proof를 함께 제시해야 한다. 이것이 Lean에서
    ``$\cdots$를 만족하는 $\cdots$''를 표현하는 표준 방식이다.
  \item \lean{def} --- definition. \lean{theorem}, \lean{lemma} --- theorem,
    lemma.
  \item \lean{Prop} --- proposition (참/거짓의 값이 아니라 \emph{proposition
    그 자체}). \lean{Bool} --- 계산 가능한 boolean 값. 이 프로젝트에서 proof
    쪽 definition은 대체로 \lean{Prop}을, 실행 가능한 algorithm 쪽 definition은
    \lean{Bool}을 쓴다. 이 구분이 \leanfile{Lean\_Algorithm.tex}의 핵심
    주제다.
  \item \lean{Option X} --- $X \cup \{\bot\}$. 값이 \lean{some v}이거나
    ``없음''을 뜻하는 \lean{none}이다.
  \item \lean{noncomputable} --- axiom of choice 등을 써서 정의했기 때문에
    실제로 계산할 수는 없고 수학적으로만 존재하는 대상이라는 표시.
  \item \lean{Set X}는 $X$의 subset (characteristic function
    $X \to \mathrm{Prop}$), \lean{Finset X}는 finite subset의 \emph{데이터}다.
    후자는 실제로 원소를 enumerate할 수 있다.
  \item \lean{Fin n} --- $\{0, 1, \dots, n-1\}$.
  \item 이름 앞의 \lean{@[simp]}는 automation tactic이 쓰는 rewrite rule로
    등록한다는 뜻이며, 수학적 내용과는 무관하다. 이하에서는 언급하지 않는다.
\end{itemize}

\begin{remark}[structure의 field는 parameter가 아니다]
\label{basics:rem:bundled}
Lean에서 자주 등장하는 설계상의 갈림길 하나를 미리 짚어 둔다. 어떤 대상의
underlying type을 \emph{parameter}로 둘 수도 있고 (예: \lean{SimpleGraph V}는
vertex type \lean{V}를 밖에서 받는다), structure의 \emph{field}로 묶어 넣을
수도 있다 (예: 아래의 \lean{τPGraph}는 vertex type \lean{V}를 자기 안에
갖고 있다). 후자를 \emph{bundled} 방식이라 한다. bundled 방식을 쓰면
``$\tau_P$-graph 하나''를 하나의 값으로 다룰 수 있어서, graph를 다른 graph로
보내는 transformation (예: incidence structure로의 translation)을 function으로
쓰기 편해진다. 이 프로젝트는 source graph 쪽은 parameter 방식 (mathlib을
그대로 씀), translation target 쪽은 bundled 방식을 쓴다.
\end{remark}

\bigskip

이하 각 절은 Lean 파일 하나(또는 서로 밀접한 몇 개)에 대응하며, 파일들이
서로를 \lean{import}하는 순서대로 배열했다.

\section{Graph 위의 \texorpdfstring{$\MSOtwo$}{MSO2} syntax}
\label{basics:sec:syntax}

\emph{파일: \leanfile{GraphMSO/Syntax.lean}}

이 파일은 이 프로젝트의 \emph{source language}, 즉 simple graph 위의
$\MSOtwo$ formula의 syntax를 정의한다. $\MSOtwo$란 vertex뿐 아니라
\emph{edge}에 대해서도 first-order variable과 (monadic) second-order variable을
허용하는 logic이다.

\subsection{Variables}

\begin{definition}[variable의 종류]
\label{basics:def:vars}
네 종류의 variable을 모두 natural number로 이름 붙인다.
\[
  \lean{FOVar} = \lean{SOVar} = \lean{EdgeFOVar} = \lean{EdgeSOVar} = \mathbb{N}.
\]
각각 vertex variable $x$, vertex-set variable $X$, edge variable $e$,
edge-set variable $E$를 뜻한다.
\end{definition}

\begin{remark}[named variable 대 de Bruijn index]
네 종류가 모두 $\mathbb{N}$이지만 \emph{서로 다른 네 개의 namespace}다.
즉 vertex variable $3$번과 edge variable $3$번은 아무 관계가 없다. 또한 이
프로젝트는 de Bruijn index가 아니라 named variable을 쓴다. 따라서 quantifier는
assignment를 덮어쓰는 방식으로 작동하며, 자동적인 capture avoidance 장치는
없다. 대신 derived formula를 정의할 때마다 bound variable을 \emph{구체적인
값}으로 고정하고, 그 선택이 옳다는 것을 해당 semantics lemma로 그때그때
확인한다. 이 방식은 \S\ref{basics:sec:treesyntax}의 tree language에서 특히
두드러진다.
\end{remark}

\subsection{Formulas}

\begin{definition}[$\MSOtwo$ formula]
\label{basics:def:formula}
\lean{Formula}는 다음 constructor로 inductive하게 정의된다.

Atomic formula:
\[
\begin{array}{ll}
  \lean{false\_} & \bot \\
  \lean{equal}\ x\ y & x = y \\
  \lean{edge}\ x\ y & x \sim y \quad (\text{vertex } x,y \text{ 가 adjacent}) \\
  \lean{inSet}\ x\ X & x \in X \\
  \lean{equalEdge}\ e_1\ e_2 & e_1 = e_2 \\
  \lean{inc}\ x\ e & \text{vertex } x \text{ 가 edge } e \text{ 에 incident} \\
  \lean{inEdgeSet}\ e\ E & e \in E
\end{array}
\]

Connective: \lean{neg}, \lean{conj}, \lean{disj}, \lean{impl},
\lean{biimpl} ($\lnot, \land, \lor, \to, \leftrightarrow$).

Quantifier: 네 종류 variable 각각에 대해 existential/universal quantifier
\lean{existsFO}, \lean{forallFO}, \lean{existsSO}, \lean{forallSO},
\lean{existsEdgeFO}, \lean{forallEdgeFO}, \lean{existsEdgeSO},
\lean{forallEdgeSO}.
\end{definition}

\begin{remark}[connective를 모두 primitive로 두는 이유]
$\lor, \to, \leftrightarrow, \forall$는 원리적으로 $\lnot, \land, \exists$로
정의할 수 있지만, 여기서는 모두 primitive constructor로 둔다. 이 덕분에 뒤의
translation (\leanfile{Lean\_MSO.tex})에서 universal quantifier를
$\lnot \exists \lnot$로 우회하지 않고 \emph{직접} guarded universal block으로
translate할 수 있고, formula의 size (\S\ref{basics:sec:cost} 참조)도 실제
구조를 반영한다.
\end{remark}

\begin{definition}[derived formula]
\label{basics:def:derived}
\lean{true\_} $:= \lnot\bot$, \lean{notEqual} $x\,y := \lnot(x=y)$,
\[
  \lean{subset}\ X\ Y := \forall v_0\,(v_0 \in X \to v_0 \in Y),
  \qquad
  \lean{setEq}\ X\ Y := (X \subseteq Y) \land (Y \subseteq X).
\]
또 variable list를 받아 차례로 quantify하는 \lean{existsFOs},
\lean{forallFOs}가 있다.
\end{definition}

\begin{remark}[bound variable $0$의 위험성]
\label{basics:rem:subsetcapture}
\lean{subset X Y}는 bound vertex variable로 \emph{항상 $0$번}을 쓴다. 따라서
이 formula를 다른 formula 안에 넣을 때 바깥에 vertex variable $0$이 free
variable로 살아 있으면 capture가 일어난다. 이 프로젝트는 이 문제를 일반론으로
해결하지 않고, \lean{subset}을 쓰는 자리마다 그 자리의 free variable이 $0$이
아님을 확인하는 방식으로 처리한다. 코드를 읽을 때 ``bound variable을 왜 굳이
$x+1$ 같은 값으로 잡았을까?''라는 의문이 들면 거의 항상 이 이유다.
\end{remark}

\subsection{Free variables and sentences}

\begin{definition}[free variable]
\label{basics:def:free}
\lean{FreeFO}, \lean{FreeSO}, \lean{FreeEdgeFO}, \lean{FreeEdgeSO}는 각각
formula $\varphi$와 variable $x$를 받아 ``$x$가 $\varphi$에서 free하게
나타난다''는 \emph{proposition}을 준다. structural induction으로 정의되며,
quantifier 경우에만
\[
  \lean{FreeFO}\ (\exists y\, \varphi)\ x \;:=\; x \neq y \;\land\; \lean{FreeFO}\ \varphi\ x
\]
처럼 bound variable을 제외한다.
\end{definition}

여기서 free variable을 \emph{set}이 아니라 \emph{proposition으로 주어지는
predicate} ($\lean{Formula} \to \lean{FOVar} \to \lean{Prop}$)로 정의한 점이
눈에 띈다. 같은 파일의 $\tau_P$ 판 (\S\ref{basics:sec:taump})에서는
\lean{Set FOVar}를 반환하도록 되어 있어 표기가 다르지만 내용은 같다.

\begin{definition}[sentence (closed formula)]
\label{basics:def:closed}
$\varphi$가 \lean{Closed}라는 것은 네 종류의 free variable이 모두 없다는 뜻,
즉 네 개의 universal proposition의 conjunction이다.
\end{definition}

\section{Graph 위의 \texorpdfstring{$\MSOtwo$}{MSO2} semantics}
\label{basics:sec:semantics}

\emph{파일: \leanfile{GraphMSO/Semantics.lean}}

\subsection{Assignment}

\begin{definition}[\lean{Assignment}]
\label{basics:def:assignment}
type $V$ (vertex)와 $E$ (edge)에 대해 \lean{Assignment V E}는 네 개의
function으로 이루어진 structure다.
\[
\begin{array}{ll}
  \lean{fo} : \mathbb{N} \to \lean{Option}\ V & \text{vertex variable} \\
  \lean{so} : \mathbb{N} \to \lean{Set}\ V & \text{vertex-set variable} \\
  \lean{efo} : \mathbb{N} \to \lean{Option}\ E & \text{edge variable} \\
  \lean{eso} : \mathbb{N} \to \lean{Set}\ E & \text{edge-set variable}
\end{array}
\]
\end{definition}

\begin{remark}[왜 first-order variable만 \lean{Option}인가]
\label{basics:rem:option}
수학책에서 assignment는 보통 ``free variable 전체에서 정의된 function''
$\rho : \mathrm{Var} \to V$로 쓴다. 그런데 Lean에서 그렇게 하면 $V$가
\emph{empty}일 때 assignment 자체가 존재하지 않게 되어, vertex가 없는
graph에서 sentence의 truth value를 논할 수 없다. 그래서 first-order variable은
``값이 없을 수도 있다''는 뜻의 \lean{Option}으로 두었다. 이 덕분에
\lean{Assignment.empty} (모든 first-order variable은 \lean{none}, 모든
second-order variable은 $\emptyset$)이 항상 존재한다. Second-order variable은
$\emptyset$이라는 자연스러운 default가 있으므로 \lean{Option}이 필요 없다.

그 대가로 \emph{atomic formula의 semantics가 통상적인 것과 미묘하게 다르다}.
아래 \ref{basics:rem:atomsem}을 보라.
\end{remark}

Assignment를 한 variable에서만 update하는 연산 \lean{updateFO},
\lean{updateSO}, \lean{updateEdgeFO}, \lean{updateEdgeSO}가 정의되고,
``update한 자리에서는 새 값, 다른 자리에서는 원래 값'', ``종류가 다른
update끼리는 commute한다'' 같은 기본 성질들이 theorem으로 붙어 있다.

\begin{definition}[\lean{AgreeOnFree}]
\label{basics:def:agree}
$\rho$와 $\sigma$가 $\varphi$에 대해 \lean{AgreeOnFree}라는 것은, $\varphi$의
free variable마다 두 assignment가 같은 값을 준다는 뜻이다 (네 종류 각각).
\end{definition}

\subsection{Satisfaction}

\begin{definition}[\lean{SatisfiesAt}]
\label{basics:def:satisfiesat}
graph $G : \lean{SimpleGraph}\ V$와 assignment
$\rho : \lean{Assignment}\ V\ G.\lean{edgeSet}$에 대한 Tarski semantics가
structural induction으로 정의된다. Connective와 quantifier는 통상적이고,
atomic formula는 다음과 같다.
\[
\begin{array}{lcl}
  G, \rho \models x = y &\iff& \exists v \in V.\ \rho(x) = \lean{some}\ v \ \land\ \rho(y) = \lean{some}\ v \\
  G, \rho \models \lean{edge}\ x\ y &\iff& \exists u, v.\ \rho(x) = \lean{some}\ u \land \rho(y) = \lean{some}\ v \land u \sim_G v \\
  G, \rho \models x \in X &\iff& \exists v.\ \rho(x) = \lean{some}\ v \land v \in \rho(X) \\
  G, \rho \models \lean{inc}\ x\ e &\iff& \exists v, \varepsilon.\ \rho(x) = \lean{some}\ v \land \rho(e) = \lean{some}\ \varepsilon \land v \in \varepsilon
\end{array}
\]
Quantifier는 \lean{updateFO} 등으로 assignment를 덮어쓴 뒤 body를 재귀적으로
평가한다.
\end{definition}

\begin{remark}[unassigned variable은 atomic formula를 false로 만든다]
\label{basics:rem:atomsem}
위 definition에서 $x$가 assign되어 있지 않으면 (\lean{none}) atomic formula는
\emph{false}다. 특히 \lean{equal x x}조차 $x$가 unassigned이면 false가 된다.
통상적인 semantics와 다른 점이지만, free variable이 있는 formula를 unassigned
assignment에서 평가할 때만 생기는 차이이므로 sentence에는 영향이 없다.
이 프로젝트의 모든 주요 theorem은 sentence (또는 free variable이 모두 assign된
assignment)에 대한 것이므로 문제되지 않는다.
\end{remark}

\begin{remark}[edge variable의 value type이 graph에 의존한다]
\label{basics:rem:edgesort}
$\rho$의 type이 \lean{Assignment V G.edgeSet}임에 주의하라. Edge variable은
``$G$의 edge 전체''라는 \emph{$G$에 의존하는 type}의 값을 갖는다. 즉
assignment의 type 자체가 graph를 언급한다. 이것은 ``edge는 항상 실제로
존재하는 edge여야 한다''는 조건을 type 수준에서 강제하는 방식이며, 그래서
$\exists e\, \varphi$의 semantics가 ``$\varepsilon \in E(G)$가 존재하여''로
자동으로 올바르게 나온다. 대신 서로 다른 graph의 assignment는 type이 달라
직접 비교할 수 없다.
\end{remark}

\begin{definition}[sentence의 satisfaction]
\label{basics:def:satisfies}
\[
  \lean{Satisfies}\ G\ \varphi \;:=\; \varphi.\lean{Closed} \ \land\
  \lean{SatisfiesAt}\ \varphi\ G\ \lean{Assignment.empty}.
\]
\end{definition}

\begin{remark}[\lean{Closed}가 definition 안에 들어 있다]
$G \models \varphi$가 ``$\varphi$가 sentence이고 \emph{또한} empty assignment
에서 satisfy된다''로 정의된 점이 특이하다. 즉 \lean{Satisfies}는 satisfaction
relation이면서 동시에 sentence임을 주장한다. 덕분에 $G \models \varphi$를
가정하는 theorem에서는 closedness 조건을 따로 가정할 필요가 없다. 반대로
$\varphi$가 sentence가 아니면 \lean{Satisfies}는 무조건 false다.
\end{remark}

\subsection{Free variable dependence}

\begin{theorem}[\lean{satisfiesAt\_ext\_on\_free}]
\label{basics:thm:extonfree}
$\rho$와 $\sigma$가 $\varphi$의 모든 free variable에서 일치하면
$G,\rho \models \varphi \iff G,\sigma \models \varphi$.
\end{theorem}

이것이 이 파일의 main theorem이며, formula에 대한 structural induction으로
증명된다. Quantifier 경우에 \lean{AgreeOnFree.updateFO} 등의 lemma를 쓴다.
Corollary로:

\begin{corollary}[\lean{satisfies\_iff\_satisfiesAt}]
\label{basics:cor:closedindep}
$\varphi$가 sentence이면, $G \models \varphi$는 \emph{임의의} assignment
$\rho$에 대해 $G, \rho \models \varphi$와 equivalent하다.
\end{corollary}

이 corollary는 뒤에서 대단히 자주 쓰인다. Sentence의 truth value를 계산할 때
아무 assignment에서나 시작해도 된다는 뜻이기 때문이다.

\section{Example formula library}
\label{basics:sec:examples}

\emph{파일: \leanfile{GraphMSO/Examples.lean}}

이 파일은 pipeline의 일부가 아니라 \emph{syntax와 semantics가 의도대로
작동하는지 확인하는 test bed}다. 전형적인 graph property를 두 번 정의한다:
한 번은 mathlib 수준의 수학적 definition으로, 한 번은
\S\ref{basics:sec:syntax}의 \lean{Formula}로. 그리고 둘이 일치한다는
theorem (\lean{satisfiesAt\_$\cdots$\_iff} 꼴)을 증명한다.

다루는 property는 다음과 같다: clique (\lean{clique}), independent set
(\lean{independent}), dominating set (\lean{dominating}), vertex cover,
bipartiteness, connectivity (\lean{connected}/\lean{disconnected}),
$k$-coloring (\lean{kColorable}, \lean{threeColorable}), perfect matching,
Hamiltonian cycle (\lean{hasHamiltonianCycleByEdges}), 그리고 $K_3$-minor
(\lean{k3Minor}, \lean{completeGraphMinor}).

\begin{remark}[connectivity와 minor를 표현하는 방식]
\label{basics:rem:connexample}
Connectivity는 ``vertex set을 두 개의 nonempty part로 나누는 어떤 partition도
그 사이를 잇는 edge를 갖는다''는 second-order formula
(\lean{connected}, \lean{isConnectedByPartition})로 표현된다. 이것이
$\MSO$에서 connectivity를 쓰는 표준적인 방법이며,
\S\ref{basics:sec:conn}에서 이 equivalence를 일반적인 형태로 증명한다.
마찬가지로 Hamiltonian cycle과 minor는 edge-set/vertex-set variable을 명시적으로
guess하는 $\MSOtwo$ formula로 표현되는데, 이것이 $\MSOtwo$가 $\MSOone$보다
strict하게 강한 이유를 보여주는 예다.
\end{remark}

\begin{remark}
Coloring이나 minor처럼 ``set 여러 개''를 다루는 formula는 \lean{List SOVar}를
받아 재귀적으로 조립된다 (\lean{coloring}, \lean{minorModelUsing},
\lean{existsSOs}). Formula의 개수가 parameter $k$에 따라 달라지는 경우를
Lean에서 다루는 표준적인 방법이다. 같은 기법이 \leanfile{Lean\_MSO.tex}의
block quantifier에서 다시 쓰인다.
\end{remark}

\section{\texorpdfstring{$\tau_P$}{tau P}-graph와 그 위의 \texorpdfstring{$\MSOone$}{MSO1}}
\label{basics:sec:taump}

\emph{파일: \leanfile{GraphMSO/language/tau\_graph.lean},
\leanfile{GraphMSO/language/syntax.lean},
\leanfile{GraphMSO/language/semantics.lean}}

\emph{참조: \leanfile{Courcelle/lecture\_note\_expanded.tex}, Theorem
\lean{thm:courcelle} (Courcelle's theorem for unary-expanded graphs)}

Lecture note의 main theorem은 vocabulary
\[
  \tau_P = \{\mathrm{adj}\} \cup P
\]
위에서 서술된다. 즉 binary adjacency relation 하나와 $P$로 index된 unary
predicate들의 family다. 보통 graph는 $P = \emptyset$인 경우다.

\begin{definition}[$\tau_P$-graph]
\label{basics:def:taupgraph}
\lean{τPGraph P}는 세 개의 field를 갖는 structure다:
vertex type \lean{V}, 그 위의 simple graph \lean{G : SimpleGraph V},
그리고 predicate interpretation \lean{pred : P → V → Prop}.
\end{definition}

\begin{remark}[$\mathrm{adj}$는 $P$의 원소가 아니다]
$\tau_P = \{\mathrm{adj}\} \cup P$라고 썼지만, Lean에서는 $\mathrm{adj}$를
$P$ 안에 넣지 않고 mandatory field \lean{G}가 interpret하도록 했다. 그래야
mathlib의 \lean{SimpleGraph} API (symmetry, irreflexivity, walk, connectivity
등)를 그대로 쓸 수 있기 때문이다. $P$는 \emph{unary predicate symbol만}
index한다.
\end{remark}

\lean{ofSimpleGraph}는 보통 graph를 $P = \lean{Empty}$인 $\tau_P$-graph로
보는 function이다. Predicate interpretation은 \lean{Empty.elim}, 즉
``empty type에서 오는 function은 유일하게 존재한다''는 사실로 준다.

\begin{definition}[$\tau_P$ 위의 $\MSOone$ syntax]
\label{basics:def:tauformula}
\lean{Language.Formula P}는 \S\ref{basics:def:formula}에서 edge 관련
constructor를 모두 빼고, 대신 predicate atom을 넣은 것이다. Atomic formula는
\lean{false\_}, \lean{equal} $x\ y$, \lean{adj} $x\ y$,
\lean{pred} $p\ x$ ($p \in P$), \lean{inSet} $x\ X$.
Connective와 vertex quantifier는 동일하다. Edge variable은 \emph{없다}.
\end{definition}

Free variable \lean{freeFO}, \lean{freeSO}는 여기서는 \lean{Set FOVar}를
반환하는 function으로 정의된다 (\S\ref{basics:def:free}의 predicate 방식과
표기만 다름).

Semantics (\lean{Language.Semantics.SatisfiesAt})는
\S\ref{basics:sec:semantics}와 완전히 평행하다: assignment는
\lean{fo : FOVar → Option X.V}와 \lean{so : SOVar → Set X.V}의 쌍이고,
predicate atom은
\[
  X, \rho \models \lean{pred}\ p\ x \iff \exists v.\ \rho(x) = \lean{some}\ v \land X.\lean{pred}\ p\ v
\]
로 해석된다. Sentence용 \lean{Language.Semantics.Satisfies}는 empty assignment
에서의 satisfaction으로 정의된다 (여기서는 \S\ref{basics:def:satisfies}와 달리
\lean{Closed} 조건을 안에 넣지 않았다).

\begin{remark}[왜 language가 두 개인가]
\label{basics:rem:twolanguages}
\S\ref{basics:sec:syntax}의 $\MSOtwo$가 우리가 \emph{묻고 싶은} language이고,
$\tau_P$ 위의 $\MSOone$이 Courcelle's theorem을 \emph{실제로 증명하는}
language다. 둘을 잇는 다리가 incidence structure
(\S\ref{basics:sec:incidence})를 거치는 translation이며, 이것이
\leanfile{Lean\_MSO.tex}의 \lean{Formula.toIncidence}다.
\end{remark}

\section{Incidence structure}
\label{basics:sec:incidence}

\emph{파일: \leanfile{GraphMSO/Decomp/KRootedPGraph.lean},
\leanfile{GraphMSO/incidence.lean}}

$\MSOtwo$를 $\MSOone$으로 reduce하는 고전적인 방법은 edge를 \emph{vertex로
승격}시키는 것이다.

\begin{definition}[\lean{KRootedPGraph}]
\label{basics:def:krootedpgraph}
\lean{KRootedPGraph P}는 \lean{τPGraph P}와 같은 데이터 (vertex type, graph,
unary predicate family)를 갖는 structure다. 이름이 다른 이유는 역할이 다르기
때문이다: 이쪽은 encoding (\leanfile{Lean\_MSO.tex})에서 rooted colored graph,
$\Sigma$-letter, incidence structure가 공유하는 \emph{공통 core}로 쓰인다.
$P$의 finiteness나 decidability는 structure에 넣지 않고 필요한 lemma에서만
가정한다.
\end{definition}

\begin{definition}[incidence graph]
\label{basics:def:incidencegraph}
$G : \lean{SimpleGraph}\ V$에 대해 \lean{IncidenceVertex G}는 두 constructor
\lean{fromV} $v$ ($v \in V$)와 \lean{fromEdge} $e$ ($e \in E(G)$)를 갖는
inductive type, 즉 $V \sqcup E(G)$다. \lean{IncidenceGraph G}는 그 위의
simple graph로, adjacency는
\[
  \lean{fromV}\ v \sim \lean{fromEdge}\ e \iff v \in e
\]
(및 그 symmetric한 경우)이고 나머지 조합은 모두 non-adjacent다. 즉 원래
vertex끼리도, edge object끼리도 adjacent하지 않은 bipartite graph다.
\end{definition}

\begin{definition}[두 개의 sort predicate]
\label{basics:def:incsort}
\lean{IncidenceVertex.IsVertex}와 \lean{IncidenceVertex.IsEdgeObj}는 각각
``원래 vertex에서 왔다'', ``edge에서 왔다''를 뜻하는 predicate다. Constructor를
보고 바로 판정되므로 \lean{Decidable} instance가 붙는다. 두 predicate가 서로의
negation이며 (\lean{isVertex\_iff\_not\_isEdgeObj}) 전체를 cover한다
(\lean{isVertex\_or\_isEdgeObj})는 것도 증명되어 있다.

\lean{IncSort}는 이 두 predicate symbol의 type
$\{\lean{vert}, \lean{edgeObj}\}$이고, \lean{colouredIncidence G}는 $G$의
incidence structure를 \lean{KRootedPGraph IncSort}로 포장한 것이다.
Vocabulary로 쓰면
$\tau_I = \{\mathrm{adj}, \mathrm{Vert}, \mathrm{EdgeObj}\}$.
\end{definition}

\begin{remark}[sort predicate가 왜 필수인가]
\label{basics:rem:whysorts}
파일의 comment가 명시하듯, vertex의 sort는 \emph{constructor에서 읽어내지
adjacency에서 복원하지 않는다}. 이것이 결정적인데, isolated된 원래 vertex와
edge object는 adjacency만 보아서는 구별되지 않기 때문이다 (둘 다 neighbor가
없거나 있을 수 있다). Sort predicate를 vocabulary에 명시적으로 넣어야만
$\MSOtwo \to \MSOone$ translation이 truth를 보존한다. 이 때문에
$P = \lean{IncSort}$인 $\tau_P$-theory가 필요한 것이다.
\end{remark}

\section{Tree decomposition의 기본 API}
\label{basics:sec:treedecomp}

\emph{파일: \leanfile{GraphMSO/Decomp/tree\_decomp.lean}}

\emph{참조: \leanfile{Courcelle/nice\_tree\_decomp.tex} \S Statement의
tree decomposition definition}

\begin{definition}[tree decomposition]
\label{basics:def:treedecomp}
Finite simple graph $G$ 위의 \lean{TreeDecomposition G}는 다음 field를 갖는
structure다.
\begin{itemize}
  \item \lean{Node}: decomposition tree의 node type, 그리고 그것이 finite하다는
    증거 \lean{nodeFintype};
  \item \lean{T : SimpleGraph Node}와 그것이 tree라는 증거 \lean{T\_istree};
  \item \lean{node2bag : Node → Set V}: 각 node의 bag;
  \item 세 개의 axiom:
  \begin{itemize}
    \item \lean{VertexCoverage}: 모든 vertex가 어떤 bag에 들어간다;
    \item \lean{EdgeCoverage}: 모든 edge의 양 endpoint가 함께 들어가는 bag이
      있다;
    \item \lean{Connectivity}: 고정된 vertex $v$를 담는 bag들의 node set이
      decomposition tree에서 preconnected하다 (running intersection).
  \end{itemize}
\end{itemize}
\lean{RootedTreeDecomposition G}는 여기에 root \lean{root : Node}를 추가한
것이다.
\end{definition}

\begin{remark}[axiom이 structure 안에 있다]
수학에서는 ``tree decomposition이란 다음 세 조건을 만족하는 $(T, \{B_t\})$
이다''라고 쓰지만, Lean에서는 그 세 조건의 \emph{proof}가 structure의 field다
(\ref{basics:rem:bundled} 참조). 따라서 \lean{TreeDecomposition G}의 값
하나를 만든다는 것은 곧 세 axiom을 증명하는 일이다. 반대로 그런 값을 갖고
있으면 \lean{D.VertexCoverage} 등으로 axiom을 즉시 꺼내 쓸 수 있다.

\lean{Connectivity}가 ``connected''가 아니라 ``preconnected'' (empty일 수도
있는 connectivity)로 되어 있는 것도 의도적이다. Nonemptiness는
\lean{VertexCoverage}에서 따라 나오므로 axiom으로 중복해 넣지 않는다.
\end{remark}

\begin{definition}[width]
\label{basics:def:haswidth}
\lean{HasWidth D ω} $:= \forall t.\ |B_t| \le \omega + 1$.
($|\cdot|$은 \lean{Set.ncard}.) 즉 width $\omega$란 bag size가 $\omega+1$
이하라는 통상적인 convention이다.
\end{definition}

\subsection{Rooted structure: depth, parent, ancestor}

Rooted tree decomposition에서 parent/child relation을 만드는 방식이 다소
우회적이므로 설명한다. Tree이므로 임의의 두 node 사이에 unique simple path가
있다는 사실 (\lean{T\_istree.existsUnique\_path})에서 출발한다.

\begin{definition}[root path, depth, parent]
\label{basics:def:parent}
\begin{itemize}
  \item \lean{rootPath T t}: root에서 $t$까지의 (unique) simple path.
    Uniqueness에서 choose하므로 \lean{noncomputable}.
  \item \lean{rootDepth T t} $:= d_T(\lean{root}, t)$, tree에서의 distance.
  \item \lean{parent T t} $:=$ \lean{rootPath T t}의 \emph{penultimate} node.
\end{itemize}
\end{definition}

\begin{remark}[root의 parent는 자기 자신]
\label{basics:rem:parentroot}
\lean{parent}는 \emph{total function}으로 정의되어야 하므로 root에서도 값을
가져야 한다. Root의 root path는 length $0$인 path이고 그 penultimate는 자기
자신이므로, \lean{parent root = root}이다. 따라서 ``$t$의 parent''라는 말을
쓸 때는 항상 $t \neq \lean{root}$를 함께 확인해야 한다. 이 때문에 진짜
parent-child relation은 별도로 정의된다:
\[
  \lean{IsChild}\ p\ c \;:=\; c \neq \lean{root} \ \land\ p = \lean{parent}\ c.
\]
이 definition 덕분에 \lean{forall\_not\_isChild\_iff} (``parent가 없는 node는
정확히 root'')가 성립하고, 이것이 tree language의 \lean{root\_} formula
(\S\ref{basics:sec:treesyntax})와 정확히 맞물린다.
\end{remark}

주요 structural lemma들:
\begin{itemize}
  \item \lean{rootDepth\_eq\_parent\_add\_one}: non-root node의 depth는
    parent보다 정확히 $1$ 크다.
  \item \lean{isChild\_of\_adj\_of\_rootDepth\_eq\_add\_one}: depth를 $1$
    늘리는 edge는 정확히 child edge다. 그 converse도 성립하므로,
  \item \lean{adj\_iff\_isChild\_or\_isChild}: tree adjacency는 두 orientation
    중 하나의 child relation이다. \emph{이 theorem이 unrooted tree와 rooted
    tree를 잇는 다리이며, tree language의 \lean{adjTree} formula가
    \lean{parent}만으로 정의될 수 있는 이유다.}
  \item \lean{IsAncestor} $:=$ \lean{IsChild}의 reflexive transitive closure
    (\lean{Relation.ReflTransGen}). 즉 ``$a$가 $b$의 ancestor''를 유한 번의
    child step으로 reachable함으로 정의한다.
  \item \lean{IsAncestor.comparable\_of\_common\_descendant}: common descendant를
    갖는 두 ancestor는 comparable하다. 증명은 ``child relation을 뒤집으면
    right-unique다'' (각 node의 parent는 unique)를 이용해 mathlib의
    \lean{Relation.ReflTransGen.total\_of\_right\_unique}를 적용한다.
\end{itemize}

\begin{definition}[subtree, cone, adhesion]
\label{basics:def:cone}
\begin{itemize}
  \item \lean{SubtreeNodes T t} $:= \{x : \lean{IsAncestor}\ t\ x\}$, $t$를
    root로 하는 subtree의 node set.
  \item \lean{cone T t} $:= \bigcup_{x \in \lean{SubtreeNodes}\ t} B_x$,
    그 subtree에 나타나는 vertex 전체 (lecture note의 \emph{cone}).
  \item \lean{adhesion T t} $:= B_t \cap B_{\lean{parent}(t)}$
    (root에서는 $\emptyset$). Relational 판인 \lean{IsAdhesion}도 함께 있다.
\end{itemize}
\end{definition}

\section{\texorpdfstring{$\mathrm{BAGS}(v)$}{BAGS(v)}와 top node}
\label{basics:sec:bags}

\emph{파일: \leanfile{GraphMSO/Decomp/bags.lean}}

\begin{definition}[$\mathrm{BAGS}(v)$]
\label{basics:def:bags}
\lean{BAGS D v} $:= \{t : v \in B_t\}$, 즉 $v$를 담는 bag들의 node set.
\end{definition}

Tree decomposition의 세 번째 axiom이 정확히 ``$\mathrm{BAGS}(v)$가 induce하는
subgraph가 preconnected하다''는 것이므로, 곧바로 \lean{BAGS\_connected}
(nonempty이므로 connected)와 \lean{BAGSGraph\_isTree} (tree의 induced subgraph는
acyclic이므로 tree)가 따라온다.

\subsection{Top node의 일반론}

이 파일의 설계에서 눈여겨볼 점은, top node 이론을 $\mathrm{BAGS}(v)$에
대해서가 아니라 \emph{임의의 preconnected node set $U$}에 대해 전개했다는
것이다.

\begin{definition}[top node]
\label{basics:def:topnode}
\lean{IsTopNode T U t} $:= t \in U \land \forall u \in U.\ \lean{rootDepth}\ t \le \lean{rootDepth}\ u$,
즉 $U$ 안에서 depth가 minimal인 원소. \lean{topNode T U hU}는 \lean{argminOn}
으로 고른 구체적인 top node다 (\lean{noncomputable}).
\end{definition}

\begin{theorem}[\lean{IsTopNode.unique}]
\label{basics:thm:topunique}
$U$가 preconnected이면 top node는 unique하다.
\end{theorem}

\begin{remark}[증명 방식]
두 top node $t \ne u$를 잡고, $U$가 induce하는 subtree 안에서 $t$부터 $u$까지의
simple path를 취한다. 그 path의 첫 step은 $t$의 minimality 때문에 depth를
\emph{늘리는} 방향이어야 하고, 그 다음부터는 simple path가 뒤로 돌아올 수
없으므로 depth가 계속 증가한다
(\lean{rootDepth\_lt\_of\_cons\_isPath\_of\_rootDepth\_eq\_add\_one}).
따라서 $\lean{rootDepth}(t) < \lean{rootDepth}(u)$인데, 이는 $u$의 minimality에
모순이다. Induced subgraph의 path를 원래 tree의 path로 옮기기 위해
injective graph homomorphism \lean{inclHom}을 따라 밀어내는 단계가 들어간다.
\end{remark}

이어지는 theorem들:
\begin{itemize}
  \item \lean{topNode\_isAncestor}: top node는 $U$의 모든 원소의 ancestor다.
  \item \lean{parent\_mem\_of\_isChild\_of\_child\_mem}: child가 $U$에 있고
    parent가 top node 아래에 있으면 parent도 $U$에 있다.
  \item \lean{mem\_of\_isAncestor\_of\_isAncestor}: $U$는 top node에서
    시작하는 ancestor chain에 대해 \emph{convex}하다.
\end{itemize}

그리고 이들을 $U = \mathrm{BAGS}(v)$에 specialize한 판
(\lean{topBAGSNode}, \lean{topBAGSNode\_isAncestor},
\lean{mem\_BAGS\_of\_isAncestor\_of\_isAncestor} 등)이 뒤따른다.

\begin{remark}[왜 일반화했는가]
\leanfile{Lean\_MSO.tex}에서 \emph{defining pair}를 다룰 때, tree language의
formula \lean{top} (\S\ref{basics:sec:treesyntax})은 ``connected node set $Z$의
top 원소''를 말한다. 이때 $Z$는 임의의 second-order variable의 값이지 반드시
$\mathrm{BAGS}(v)$ 꼴이 아니다. 일반론을 먼저 세워 두었기 때문에 formula의
semantics와 combinatorial fact가 곧바로 맞물린다.
\end{remark}

\section{Bag coloring}
\label{basics:sec:bagcoloring}

\emph{파일: \leanfile{GraphMSO/Decomp/bagColoring.lean}}

\begin{definition}[bag-injective coloring]
\label{basics:def:bagcoloring}
\lean{BagColorSet ω} $:= \lean{Fin}(\omega+1) = \{0,\dots,\omega\}$.
Coloring $c : V \to \lean{Fin}\ k$가 \lean{IsBagColoring}이라는 것은 모든
bag $B_t$ 위에서 $c$가 injective라는 뜻 (\lean{Set.InjOn}).
\lean{BagColoring D ω}는 그런 성질을 갖는 $c : V \to \lean{Fin}(\omega+1)$들의
subtype이다.
\end{definition}

\begin{definition}[conflict graph]
\label{basics:def:conflictgraph}
\lean{bagConflictGraph T}는 $V$ 위의 graph로, $u \sim v$를 ``$u \ne v$이고
어떤 bag이 둘을 함께 담는다''로 정의한다.
\end{definition}

\begin{lemma}[\lean{isBagColoring\_iff\_valid\_conflict}]
\label{basics:lem:conflictiff}
Bag-injective coloring은 정확히 conflict graph의 proper coloring이다.
\end{lemma}

이 관점 전환이 다음 theorem의 핵심이다.

\begin{theorem}[\lean{exists\_bagColoring\_of\_hasWidth}]
\label{basics:thm:existsbagcoloring}
Width $\omega$ 이하의 rooted tree decomposition은 $\{0,\dots,\omega\}$를 쓰는
bag-injective coloring을 갖는다.
\end{theorem}

\begin{remark}[증명 방식: greedy + top node argument]
\label{basics:rem:greedyproof}
증명은 두 layer로 나뉜다.

\emph{(1) 일반적인 greedy coloring lemma}
(\lean{exists\_coloring\_of\_forall\_exists\_neighbor\_card\_le}):
finite graph $H$에서 ``nonempty인 모든 finite subset $s$가 $s$ 안에서
neighbor가 $\omega$개 이하인 vertex를 갖는다''면 $H$는
$(\omega+1)$-colorable하다. 증명은 \lean{Finset.strongInduction}으로 $s$에 대한
strong induction을 쓴다: degree가 낮은 vertex $v$를 하나 골라
$s \setminus \{v\}$를 먼저 color하고, $v$의 neighbor들이 쓴 color가 최대
$\omega$개이므로 $\omega+1$개 중 남는 color가 있다.

\emph{(2) conflict graph가 그 조건을 만족한다는 것}: finite set $s$에서
$\lean{topBAGSDepth}(v) = \lean{rootDepth}(\lean{topBAGSNode}(v))$가 maximal인
$v$를 고른다. $u$가 $s$ 안의 conflict neighbor이면
$\lean{topBAGSDepth}(u) \le \lean{topBAGSDepth}(v)$이고, 이때
\lean{mem\_bag\_topBAGSNode\_of\_conflict\_of\_depth\_le}에 의해
$u \in B_{\lean{topBAGSNode}(v)}$이다. 즉 \emph{$v$의 conflict neighbor는 모두
$v$의 top bag 안에 들어 있다}. 그 bag의 size는 $\omega+1$ 이하이고 $v$ 자신을
빼면 $\omega$ 이하이므로 (1)의 가정이 성립한다.

(2)의 핵심인 \lean{mem\_bag\_topBAGSNode\_of\_conflict\_of\_depth\_le}는
\S\ref{basics:sec:bags}의 convexity로 증명된다: $u, v$가 함께 있는 bag $t$가
있으면 $\lean{topBAGSNode}(u)$와 $\lean{topBAGSNode}(v)$는 모두 $t$의
ancestor이므로 comparable하고, depth 조건에서 전자가 후자의 ancestor다. 그러면
$\lean{topBAGSNode}(v)$는 $\mathrm{BAGS}(u)$의 top node와 원소 $t$ 사이에
끼어 있으므로 convexity에 의해 $\mathrm{BAGS}(u)$에 속한다.
\end{remark}

\section{Edge bound}
\label{basics:sec:edgebound}

\emph{파일: \leanfile{GraphMSO/Decomp/edgeBound.lean}}

\emph{참조: \leanfile{Courcelle/lecture\_note\_expanded.tex} Lemma
\lean{lem:tw-edges}}

\begin{theorem}[\lean{edgeSet\_ncard\_le\_of\_hasWidth}]
\label{basics:thm:edgebound}
$G$가 width $\omega$ 이하의 tree decomposition을 가지면
$|E(G)| \le \omega \cdot |V(G)|$.
\end{theorem}

\begin{remark}[lecture note와 다른 증명]
\label{basics:rem:edgeboundproof}
Lecture note는 minimal decomposition에 대한 induction으로 증명하지만, Lean
증명은 \emph{charging argument}를 쓴다. 각 edge $\{u,v\}$를
$\lean{topBAGSDepth}$가 더 \emph{깊은} 쪽 endpoint에 charge한다.
\S\ref{basics:rem:greedyproof}의 lemma에 의해 다른 endpoint는 charge된
vertex의 top bag 안에 있으므로, vertex $v$에 charge되는 edge는
$B_{\lean{topBAGSNode}(v)} \setminus \{v\}$의 원소 개수, 즉 $\omega$개
이하다. 따라서 총 edge 수는 $\omega \cdot |V|$ 이하다.

Formalization 관점에서 이 선택이 중요한 이유는, minimal decomposition induction이
``decomposition을 surgery해서 더 작게 만든다''는 절차를 formalize해야 하는
반면, charging argument는 \S\ref{basics:sec:bags}에서 이미 세운 top node
machinery를 재사용하기만 하면 되기 때문이다.
\end{remark}

Unrooted 판 (\lean{TreeDecomposition.edgeSet\_ncard\_le\_of\_hasWidth})은 아무
node나 root로 잡아 rooted 판에 넘기는 방식으로 얻는다.

\section{Induced connectivity의 partition characterization}
\label{basics:sec:conn}

\emph{파일: \leanfile{GraphMSO/connectivity.lean}}

\begin{theorem}[\lean{induce\_preconnected\_iff\_forall\_exists\_adj}]
\label{basics:thm:connpartition}
임의의 simple graph $G$와 $X \subseteq V$에 대해,
\[
  G[X] \text{ 가 preconnected} \iff
  \forall Y \subseteq X,\ Y \ne \emptyset,\ X \setminus Y \ne \emptyset
  \implies \exists u \in Y, v \in X \setminus Y.\ u \sim_G v.
\]
따라서 \lean{induce\_connected\_iff\_nonempty\_and\_forall\_exists\_adj}:
$G[X]$가 connected $\iff$ $X \ne \emptyset$이고 위 조건.
\end{theorem}

\begin{remark}[이 theorem의 역할과 증명]
이것이 tree language의 \lean{conn}$(X)$ formula
(\S\ref{basics:sec:treesyntax})의 \emph{semantic content}다. $\MSO$는
``path가 존재한다''를 직접 말할 수 없으므로, connectivity를 ``어떤 proper
subset도 edge에 대해 closed되어 있지 않다''는 second-order 조건으로 표현한다.
이 theorem이 그 표현이 옳음을 보장한다.

($\Rightarrow$) 방향은 $Y$의 한 점에서 $X \setminus Y$의 한 점으로 가는 walk를
잡고, 그 walk를 따라가며 처음으로 $Y$를 벗어나는 edge를 찾는다 (walk에 대한
induction). ($\Leftarrow$) 방향은 한 점 $u$에서 reachable한 점들의 set $Y$를
잡는다. $Y$를 crossing하는 edge가 있으면 그 edge의 반대편도 reachable해져
모순이므로 $X \setminus Y = \emptyset$, 즉 $Y = X$다.

Finiteness 가정이 전혀 필요 없다는 점도 기록해 둘 만하다.
\end{remark}

\section{Labeled tree 위의 \texorpdfstring{$\MSO$}{MSO} syntax}
\label{basics:sec:treesyntax}

\emph{파일: \leanfile{GraphMSO/treeLanguage/syntax.lean}}

\emph{참조: \leanfile{Courcelle/lecture\_note\_expanded.tex}
\S``Converting graph MSO to tree MSO'', 특히 \S``Basic tree formulas''
및 \S``Recognizing vertices and vertex sets''}

이것이 Courcelle translation의 \emph{target language}다. Formula는 labeled
rooted tree에 대해, parent relation과 label predicate를 통해 말한다.

\begin{definition}[tree formula]
\label{basics:def:treeformula}
Alphabet $A$에 대해 \lean{TreeLanguage.Formula A}의 atomic formula는:
\[
\begin{array}{ll}
  \lean{equal}\ x\ y & x = y \\
  \lean{parent}\ y\ x & \text{$y$가 $x$의 parent} \\
  \lean{labelMem}\ S\ x & x \text{의 label} \in S \quad (S \subseteq A) \\
  \lean{labelMem₂}\ R\ x\ y & (x \text{의 label}, y \text{의 label}) \in R \quad (R \subseteq A \times A) \\
  \lean{inSet}\ x\ X & x \in X
\end{array}
\]
Connective와 node quantifier는 이전과 같다.
\end{definition}

\begin{remark}[lecture note와의 두 가지 차이 --- 둘 다 semantics 보존적]
\label{basics:rem:treesimplify}
파일 comment가 명시하는 두 simplification을 반드시 짚어야 한다.

\emph{(1) \lean{child₁}, \lean{child₂} 대신 \lean{parent}.}
Lecture note의 vocabulary는 ordered child relation \lean{child₁},
\lean{child₂}를 갖지만, translation 부분에서는 사실상 abbreviation
$\lean{parent}(y,x) = \lean{child₁}(y,x) \lor \lean{child₂}(y,x)$로만 쓴다.
그래서 이 formalization은 \lean{parent}를 primitive로 둔다. Ordered child
relation은 tree를 automaton에 넘기는 지점에서만 필요하고, 그곳에서는
\lean{parent}가 first-order definable하다 (\leanfile{Lean\_Automata.tex} 참조).

\emph{(2) letter별 predicate $P_a$ 대신 letter \emph{set} predicate
\lean{labelMem}.}
Lecture note에서 $P_a$는 항상 finite disjunction
$\bigvee_{a \in S} P_a(x)$의 형태로만 등장한다 ($S$는 letter class
$A_i, R_i, E_{ij}, Q_i$ 중 하나). 그래서 이 formalization은 ``label이 $S$에
속한다''를 \emph{하나의 atomic formula}로 둔다. $P_a$ 자체는 $S = \{a\}$인
특수한 경우 (\lean{labelIs})다. 마찬가지로 \lean{labelMem₂ R x y}는 legality
formula에 나오는 double disjunction
$\bigvee_{(a,b) \in R} (P_a(x) \land P_b(y))$를 하나의 atomic formula로
대체한다.

이 simplification은 alphabet이 finite할 때 expressive power를 바꾸지 않으면서
formula의 size와 증명의 부담을 크게 줄인다. 대신 ``formula의 size''를 셀 때
atomic formula 하나가 실제로는 alphabet size만큼의 disjunction임을 기억해야
한다.
\end{remark}

\begin{definition}[derived tree formula]
\label{basics:def:treederived}
\begin{itemize}
  \item \lean{root\_} $x := \lnot \exists x{+}1.\ \lean{parent}(x{+}1, x)$
    --- $x$는 parent가 없다.
  \item \lean{subset}, \lean{setEq}, \lean{empty}, \lean{nonempty}
    --- set 연산.
  \item \lean{adjTree} $x\ y := \lean{parent}(x,y) \lor \lean{parent}(y,x)$
    --- undirected tree edge.
  \item \lean{conn} $X$ --- $X$가 nonempty이고, $X$의 모든 partition이 tree
    edge로 crossing된다 (\S\ref{basics:thm:connpartition}의 formula 판).
  \item \lean{top} $x\ X$ --- $x \in X$, $X$가 connected, 그리고 $x$는 root이거나
    parent가 $X$ 밖에 있다.
  \item \lean{dangle} $x\ X$ --- $x \notin X$이지만 parent는 $X$ 안에 있다.
  \item \lean{existsSOList}, \lean{forallSOList} --- set variable list에 대한
    quantifier block.
  \item \lean{conjList}, \lean{disjList} --- finite conjunction/disjunction.
\end{itemize}
\end{definition}

\begin{remark}[bound variable의 구체적 선택]
\label{basics:rem:treebound}
\ref{basics:rem:subsetcapture}에서 예고한 대로, derived formula의 bound
variable은 구체적으로 고정되어 있다: \lean{root\_ x}는 $x+1$, \lean{subset}은
vertex variable $0$, \lean{conn X}는 set variable $X+1$과 vertex variable
$0, 1$, \lean{top x X}는 $x+1$ 등. 일반적인 capture avoidance 장치를 만드는
대신, 각 formula의 semantic characterization theorem
(\S\ref{basics:sec:treesem})이 그 선택이 옳음을 그때그때 verify한다.
\end{remark}

Recognition formula들도 이 파일에서 \emph{syntax로만} 정의된다.
이들의 semantic property는 \leanfile{Lean\_MSO.tex}에서 다룬다.

\begin{definition}[legality와 recognition formula]
\label{basics:def:recognitionformulas}
\begin{itemize}
  \item \lean{legal S R} --- root의 letter가 $S$에 속하고, 모든 non-root
    node의 letter가 parent의 letter와 $R$-compatible하다.
    (Lecture note의 $\varphi_{\mathrm{legal}}$.)
  \item \lean{definingPair SA SR Z} --- lecture note의
    $\varphi_{\mathrm{vtx}\text{-}i}$. $Z$가 connected이고, $Z$의 모든 letter가
    $SA$에, top이 아닌 원소는 $SR$에 속하며, top 원소와 dangling child들은
    $SR$에 속하지 않는다. \emph{lecture note에 있는 legality conjunct는 뺐다}:
    encoding 위에서는 자동으로 성립하고, 최종 sentence에서 top level에 한 번만
    붙이기 때문이다.
  \item \lean{vtxTuple}, \lean{setTuple} --- vertex 하나 / vertex set 하나를
    recognize하는 formula (lecture note의 $\varphi_{\mathrm{vtx}}$,
    $\varphi_{\mathrm{set}}$).
  \item \lean{adjTuple}, \lean{predTuple}, \lean{eqTuple}, \lean{contTuple}
    --- atomic formula $\mathrm{adj}$, unary predicate, equality, membership을
    tuple 수준에서 표현 (lecture note의
    $\varphi_{\mathrm{adj}}, \varphi_Q, \varphi_{=}, \varphi_{\mathrm{cont}}$).
\end{itemize}
이들은 모두 color $i \in \lean{Fin}\ k$에 대한 finite disjunction/conjunction
(\lean{disjList}/\lean{conjList} + \lean{List.finRange k})으로 조립된다.
\end{definition}

\section{Tree language의 semantics}
\label{basics:sec:treesem}

\emph{파일: \leanfile{GraphMSO/treeLanguage/semantics.lean}}

\begin{definition}[tree model]
\label{basics:def:treemodel}
\lean{TreeModel A}는 세 field를 갖는 structure다: node type \lean{Node},
parent relation \lean{parentRel : Node → Node → Prop}, label
\lean{label : Node → A}.
\end{definition}

\begin{remark}[tree axiom이 \emph{없다}]
\label{basics:rem:notree}
결정적으로 중요한 설계 결정이다. \lean{TreeModel}에는 ``실제로 tree다''라는
axiom이 전혀 들어 있지 않다. Parent relation은 그냥 임의의 binary relation이다.

이유는 두 가지다. 첫째, $\MSO$ semantics 자체는 tree 성질을 필요로 하지 않는다.
둘째, tree 성질을 structure에 넣으면 모든 lemma가 그 proof를 끌고 다녀야 하고,
``legality formula가 성립하는 model''처럼 tree 성질이 \emph{conclusion}인
상황을 다루기가 어려워진다. 의도된 model은 decomposition의 $\Sigma$-tree
encoding이며, tree property가 필요한 theorem은 그때그때 명시적으로 가정하거나
encoding에서 물려받는다.
\end{remark}

\begin{definition}[symmetrized parent graph]
\label{basics:def:treemodelgraph}
\lean{TreeModel.graph M} $:= \lean{SimpleGraph.fromRel}\ M.\lean{parentRel}$,
즉 $u \sim v \iff u \ne v \land (\lean{parentRel}\ u\ v \lor \lean{parentRel}\ v\ u)$.
\end{definition}

Semantics \lean{SatisfiesAt}은 \S\ref{basics:sec:semantics}와 같은 꼴이며,
\lean{Satisfies}는 empty assignment에서의 satisfaction이다.

\subsection{Block quantifier와 \lean{setBlock}}

\begin{theorem}[\lean{satisfiesAt\_existsSOList\_iff}, \lean{satisfiesAt\_forallSOList\_iff}]
\label{basics:thm:blockquant}
Set variable list $l$에 대한 existential quantifier block이 성립한다는 것은,
``$l$ 밖의 variable에서는 $\rho$와 일치하는 어떤 assignment $\rho'$이 body를
satisfy한다''는 것과 equivalent하다. Universal 판도 마찬가지.
\end{theorem}

이것이 ``순차적으로 하나씩 quantify한다''는 syntax와 ``동시에 여러 개를
바꾼다''는 semantics를 잇는 다리다. 후자를 직접 표현하기 위해 다음이 정의된다.

\begin{definition}[\lean{Assignment.setBlock}]
\label{basics:def:setblock}
$Zs : \lean{Fin}\ k \to \lean{SOVar}$와 $g : \lean{Fin}\ k \to \lean{Set}\ \lean{Node}$에
대해, \lean{ρ.setBlock Zs g}는 $Zs$가 enumerate하는 variable들만 $g$의 값으로
덮어쓰고 나머지는 그대로 두는 assignment다.
\end{definition}

\begin{remark}[\lean{setBlock}의 definition이 까다로운 이유]
\label{basics:rem:setblock}
Lean definition은
\[
  \lean{so}\ Y := \text{if } h : \exists i.\ Zs(i) = Y \text{ then } g(h.\lean{choose}) \text{ else } \rho.\lean{so}\ Y
\]
이다. 즉 ``$Y$가 block에 속하면 그 index의 값을 준다''인데, 그 index를
\emph{axiom of choice로 뽑는다} (\lean{h.choose}). $Zs$가 injective가 아니면
어느 index가 뽑힐지 알 수 없으므로, 원하는 equation
\[
  (\rho.\lean{setBlock}\ Zs\ g).\lean{so}(Zs(i)) = g(i)
\]
는 무조건 성립하지 않는다. 그래서 이 equation (\lean{setBlock\_so\_apply})은
\emph{$Zs$의 injectivity를 hypothesis로} 갖는다. 이것이
\leanfile{Lean\_MSO.tex}의 variable allocation (각 graph variable에 $\omega+1$
개의 tree set variable block을 배정하는 부분)에서 injectivity와 disjointness
lemma를 따로 증명해야 하는 이유다.
\end{remark}

\subsection{Derived formula의 semantic characterization}

각 derived formula에 대해 ``$\models$ 인 것과 $\cdots$인 것이 equivalent''를
주장하는 theorem이 하나씩 있다: \lean{satisfiesAt\_labelMem\_iff},
\lean{\_nonempty\_iff}, \lean{\_empty\_iff}, \lean{\_subset\_iff},
\lean{\_setEq\_iff}, \lean{\_root\_iff}, \lean{\_adjTree\_iff},
\lean{\_inSet\_iff}, \lean{\_dangle\_iff}, \lean{\_labelMem₂\_iff},
\lean{\_conjList}, \lean{\_disjList}, \lean{\_legal\_iff}.
First-order variable을 언급하는 것들은 ``$\rho(x) = \lean{some}\ n$''이라는
hypothesis를 달고 있는데, 이는 \ref{basics:rem:atomsem}에서 말한 unassigned
문제를 피하기 위함이다.

두 개만 따로 언급한다.

\begin{theorem}[\lean{satisfiesAt\_conn\_iff}]
\label{basics:thm:conniff}
$M, \rho \models \lean{conn}(X) \iff$ $\rho(X)$가 symmetrized parent graph
\lean{TreeModel.graph M}에서 connected induced subgraph를 이룬다.
\end{theorem}

증명은 \S\ref{basics:thm:connpartition}으로 곧장 reduce된다. 다만 formula의
bound variable $X+1$이 $X$와 다르다는 사실 ($X \ne X+1$)을 명시적으로 써야
하고, formula의 ``crossing edge'' 조건과 \lean{fromRel}의 adjacency ($u \ne v$
조건이 추가로 붙는다) 사이의 간극을 메워야 한다. 후자는 어렵지 않다: crossing
edge의 두 endpoint는 한쪽은 $Y$ 안, 한쪽은 $Y$ 밖이므로 자동으로 서로 다르다.

\begin{theorem}[\lean{satisfiesAt\_top\_iff}]
\label{basics:thm:topiff}
$M, \rho \models \lean{top}(x, X) \iff$
$n \in \rho(X)$이고 $\rho(X)$가 connected이며,
$n$이 parent가 없거나 parent가 $\rho(X)$ 밖에 있다.
\end{theorem}

\S\ref{basics:def:topnode}의 combinatorial top node와 이 formula가 encoding
위에서 일치한다는 것이 \leanfile{Lean\_MSO.tex}의 핵심 lemma 중 하나다.

\section{Tree model의 isomorphism}
\label{basics:sec:modeliso}

\emph{파일: \leanfile{GraphMSO/treeLanguage/modelIso.lean}}

\begin{definition}[\lean{TreeModel.Iso}]
\label{basics:def:treeiso}
$M$과 $N$ 사이의 isomorphism은 node의 bijection \lean{toEquiv}와,
그것이 parent relation을 보존한다는 증거, label을 보존한다는 증거의 삼중쌍이다.
\end{definition}

\begin{theorem}[\lean{satisfiesAt\_mapEquiv\_iff}, \lean{satisfies\_iff\_of\_iso}]
\label{basics:thm:isoinvariance}
Tree $\MSO$ satisfaction은 model의 isomorphism에 대해 invariant하다.
\end{theorem}

증명은 formula에 대한 structural induction이다. Assignment를 bijection을 따라
옮기는 연산 \lean{Assignment.mapEquiv} (first-order는 \lean{Option.map},
second-order는 image)를 먼저 정의하고, 그것이 update 연산과 commute한다는
lemma (\lean{mapEquiv\_updateFO}, \lean{mapEquiv\_updateSO})를 쓴다.
Second-order quantifier 경우에는 $S$의 preimage $e^{-1}(S)$를 witness로 쓰고
$e(e^{-1}(S)) = S$ (surjectivity)를 이용한다.

\begin{remark}[이 theorem이 필요한 이유]
이것이 최종 assembly의 \emph{glue}다. Translation correctness theorem은
decomposition의 (unordered) $\Sigma$-tree encoding model에 대해 말하는데,
automaton은 ordered binary tree 위에서 돈다. 두 model 사이의 isomorphism을
만들면 (\leanfile{Lean\_Automata.tex}의 \lean{orderedEncodeIso}) 이 theorem이
satisfaction을 옮겨 준다.
\end{remark}

\section{Pendant extension}
\label{basics:sec:pendant}

\emph{파일: \leanfile{GraphMSO/pendant.lean}}

\begin{definition}[\lean{pendantExtension}]
\label{basics:def:pendant}
Graph $G$ (vertex $\alpha$)와 attachment function
$\lean{attach} : \beta \to \alpha$에 대해,
\lean{pendantExtension G attach}는 vertex $\alpha \sqcup \beta$ 위의 graph로,
$G$의 edge를 모두 유지하고 각 $b \in \beta$를 $\lean{attach}(b)$에 leaf로
붙인 것이다.
\end{definition}

주요 theorem: \lean{preconnected}, \lean{connected}, \lean{isAcyclic},
\lean{isTree} --- pendant extension은 connectivity, acyclicity, treeness를 모두
보존한다.

\begin{remark}[acyclicity 증명]
\label{basics:rem:pendantacyclic}
Connectivity는 쉽지만 acyclicity는 약간의 작업이 필요하다. 증명은 두 단계다.
먼저 pendant leaf는 degree가 $1$이므로 어떤 cycle에도 놓일 수 없음을 보인다
(\lean{IsCycle.not\_mem\_support\_of\_forall\_adj\_eq}: cycle 위의 점은
서로 다른 neighbor 두 개를 가져야 한다). 그러면 임의의 cycle은 $\alpha$ 쪽에만
놓이고, 그것을 induced subgraph $G$의 cycle로 끌어올린다
(\lean{induceLift}, \lean{isCycle\_induceLift\_iff}). $G$가 acyclic이므로
모순이다.
\end{remark}

이 파일은 \leanfile{Lean\_MSO.tex}의 incidence graph decomposition 구성에서
쓰인다: $G$의 decomposition에 각 edge마다 leaf를 하나씩 붙여
$\lean{IncidenceGraph}\ G$의 decomposition을 만들 때, 새 decomposition tree가
여전히 tree임을 이 파일이 보장한다.

\section{Abstract cost layer}
\label{basics:sec:cost}

\emph{파일: \leanfile{GraphMSO/Cost.lean}}

\begin{definition}[\lean{Costed}]
\label{basics:def:costed}
\lean{Costed α}는 value \lean{value : α}와 operation count \lean{cost : ℕ}의
쌍이다. 연산:
\[
\begin{array}{ll}
  \lean{pure}\ a = (a, 0) & \text{cost 없이 value 반환} \\
  \lean{tick}\ a = (a, 1) & \text{primitive operation 1회} \\
  \lean{charge}\ n\ a = (a, n) & \text{명시적으로 $n$회} \\
  \lean{bind}\ x\ f = (f(x.\lean{value}).\lean{value},\ x.\lean{cost} + f(x.\lean{value}).\lean{cost}) & \text{sequential composition} \\
  \lean{map}\ f\ x = (f(x.\lean{value}),\ x.\lean{cost}) & \text{pure function 적용}
\end{array}
\]
\end{definition}

\begin{remark}[무엇을 주장하고 무엇을 주장하지 않는가]
\label{basics:rem:costmeaning}
이 layer는 Lean의 실제 evaluator (kernel이나 VM)와 \emph{무관하다}. ``어떤
operation에 $1$을 매길 것인가''는 사용하는 쪽이 정하고, proof는 그 결과
counter를 compositional하게 추적할 뿐이다. 따라서
\leanfile{Lean\_Algorithm.tex}에 나오는 ``cost가 정확히 $3n+2$''류의 theorem은
\emph{abstract unit-cost model에서의 operation count}에 대한 주장이지, Lean
program의 실제 실행 시간에 대한 주장이 아니다. 이 점은
\leanfile{docs/PLAN.md}에도 명시되어 있다.

이런 layer를 두는 이유는, Courcelle's theorem의 ``linear time'' 부분을
형식적으로 말하려면 time model이 필요한데, Lean의 execution model을
formalize하는 것은 전혀 다른 규모의 작업이기 때문이다. 대신 algorithm의 구조가
강제하는 operation count를 정확히 세는 것으로 타협한다.
\end{remark}

\lean{ticks n}은 정확히 $n$번의 unit operation을 수행하며, \lean{ticks\_cost}가
그 cost가 $n$임을 증명한다.

\end{document}
